\documentclass{report}
\usepackage{amsfonts, amsmath}
\newcommand\R{{\mathbb R}}
\newcommand\Q{{\mathbb Q}}
\newcommand\N{{\mathbb N}}
\newcommand\ve{\varepsilon}
\newcommand\Co{{\mathcal C}}
\pagestyle{empty}
\begin{document}
\large
\centerline{\Huge Fisica Matematica II}
\renewcommand{\thefootnote}{ } 
\footnote{\sl Avete 2:30 ore di tempo. Ogni esercizio
vale dieci punti. La sufficienza si ottiene con un punteggio $\geq$
18. Solo le risposte chiaramente giustificate saranno prese in
considerazione.}
\renewcommand{\thefootnote}{\arabic{footnote}} 
\vskip.1cm
\centerline{\Large Appello 28-01-2004}
\vskip1cm
\begin{enumerate}
\item Sia $Q:=\{(x,y)\in\R^2\;:\; x^2+y^2\leq 1\}$ e sia $u:\bar Q\to\R$ una
soluzione, non identicamente nulla, di
\[
\begin{split}
&\Delta u=-\lambda u\\
&u|_{\partial Q}=0.
\end{split}
\]
Si dimostri che deve essere $\lambda >0$. Si cerchi di determinare il
$\lambda$ minimo possibile.
\item Si risolva l'equazione
\[
\begin{split}
&(1-y^2)u_y+u_x=u\\
&u(x,0)=x.
\end{split}
\]
\item Si determini la soluzione dell'equazione
\[
\begin{split}
&u_{tt}=u_{xx}\quad\quad 0<x<1\quad t>0\\
&u(x,0)=x\quad u_t(x,0)=0 \quad 0\leq x\leq 1\\
&u_x(0,t)=u(1,t)=0
\end{split}
\]
Si discuta la regolarit\`a della soluzione trovata.
\item Si determini la soluzione del problema
\[
\begin{split}
&u_{xx}+xe^{-t}=u_t\quad x\in[0,1],\;t>0\\
&u(x,0)=\sin 2\pi x\quad x\in[0,1]\\
&u(t,0)=u(t,1)=0\quad t>0.
\end{split}
\]
\end{enumerate}
\end{document}
